% !TeX root = ../main.tex
% ─────────────────────────────────────────────────────────────
%  NGUỒN GỐC — Cờ Trận là trò chơi dân gian; tác giả KHÔNG sáng tạo ra nó.
%  Mọi thông tin ở đây đến từ hồi ức của tác giả (14/10/2026).
% ─────────────────────────────────────────────────────────────
\section{Nguồn Gốc và Phương Pháp Ghi Chép}
\label{sec:provenance}

\subsection{Nơi chơi, người chơi và tên gọi}

Cờ Trận là trò chơi của trẻ em ở thôn Đồng Bùa, xã Tam Quan, huyện Tam
Đảo, tỉnh Vĩnh Phúc. Bài báo dùng địa danh theo địa giới hành chính vào
thời điểm trò chơi được chơi; từ năm 2025, tỉnh Vĩnh Phúc được hợp nhất
vào tỉnh Phú Thọ. Tác giả được dạy chơi Cờ Trận khi học lớp 5, khoảng năm
học 2011--2012, rồi chơi cùng bạn bè trong làng đến hết lớp 9, khoảng năm
2016, thì thôi. Như vậy trò chơi đã được truyền lại cho tác giả từ những
người chơi trước đó, chứ không bắt đầu từ thế hệ của tác giả.

Một ván 4 người thường kéo dài 3 đến 4 tiếng, từ trưa đến chiều tối. Trò
chơi dân gian \emph{không} có giới hạn số nước: ván cờ chỉ kết thúc theo
điều kiện thắng ở Mục~\ref{sec:rules}.

``Cờ Trận'' là tên người trong làng dùng. Chữ \emph{trận} gợi đến trận
đánh và thế trận. Tên tiếng Anh \emph{Tactical Formation} là cách tác giả
chuyển ngữ tên gọi này cho người đọc quốc tế; đó không phải tên dân gian.

\subsection{Bàn cờ và quân cờ}

Trò chơi không cần đồ chơi chế sẵn. Bàn cờ được kẻ tuỳ ý ngay trên đường.
Nếu là đường đất, người chơi dùng que cây vạch lên mặt đất. Nếu là đường bê
tông, người chơi kẻ bằng đá, phấn hoặc gạch để nét vẽ lâu mất hơn. Quân cờ
là bất cứ thứ gì nhặt được ven đường: viên đá, hòn sỏi, cành cây khô hoặc
tươi, chiếc lá. Mỗi phe dùng một loại vật khác nhau. Khi ăn quân, người
chơi nhặt vật của phe bị ăn ra và đặt vật của phe mình vào đúng chỗ đó. Cách
làm này giữ nguyên tổng số quân trên bàn, đúng với định luật bảo toàn quân ở
Phụ lục~\ref{app:formal}. Việc dùng vật nhặt được cũng giải thích vì sao trò
chơi không để lại hiện vật nào.

\subsection{Phương pháp ghi chép và mức độ tin cậy}

Luật chơi trong bài báo được tác giả ghi lại từ hồi ức của chính mình.
Theo tác giả, bản ghi gồm đầy đủ và chính xác các thành phần sau, đúng như
trò chơi dân gian: bàn cờ hai lớp đỉnh--tâm, cách nhảy, luật kẹp và việc
chỉ chọn một hướng ăn, số quân và cách xếp quân ban đầu, điều kiện thắng,
và cả ba chế độ 2, 3 và 4 người.

Chúng tôi chưa thể đối chiếu bản ghi này với nguồn thứ hai. Hầu hết người
trong làng đã quên trò chơi. Tìm kiếm trên web bằng tiếng Việt và tiếng Anh
(tháng 10/2026) cũng không cho thấy tài liệu nào mô tả Cờ Trận hoặc một trò
chơi có cùng bộ luật. Như vậy bản ghi dựa trên \emph{một người cung cấp
thông tin duy nhất}, và chúng tôi nêu rõ giới hạn này ở
Mục~\ref{sec:limitations}. Chúng tôi mong những ai biết Cờ Trận, hoặc một
biến thể ở làng khác, liên hệ với tác giả để bổ sung và đối chiếu.

\subsection{Phần dân gian và phần bổ sung}

Để người đọc không nhầm lẫn giữa trò chơi dân gian và đóng góp của bài báo,
chúng tôi tách rõ hai phần:
\begin{itemize}
  \item \textbf{Thuộc trò chơi dân gian}: toàn bộ luật chơi được đặc tả ở
        Mục~\ref{sec:rules} và Mục~\ref{sec:variants}.
  \item \textbf{Do tác giả bổ sung}: tên tiếng Anh; mô hình toán học hình
        thức (Phụ lục~\ref{app:formal}); phần mềm và giao diện
        (Mục~\ref{sec:impl}); thứ tự lượt cố định trong phần mềm
        (Mục~\ref{ssec:turn_order}); toàn bộ AI (Mục~\ref{sec:ai});
        và các thực nghiệm (Mục~\ref{sec:eval}).
\end{itemize}
